\documentclass[10pt,a4paper]{amsart}
\usepackage[margin=19mm]{geometry}
\usepackage{amsmath,amssymb,mathtools}
\usepackage[scr=rsfs,cal=euler]{mathalfa}
\usepackage{xfrac}
\usepackage[unicode,hidelinks,bookmarks=false]{hyperref}
\newcommand{\ie}{i.e.\ }
\newcommand{\cf}{cf.\ }
\newcommand{\resp}{respectively\ }
\newcommand{\Subsection}{\subsection}
\providecommand{\nobreakdash}{\nobreak\hbox{-}}
\setlength{\emergencystretch}{2em}
\multlinegap0pt
\usepackage{mathtools}
\usepackage{amssymb}
\usepackage{xcolor}
\newtheorem{theorem}{Theorem}
\newcommand{\AJS}[1]{#1}
\newcommand{\Real}{\mathbb{R}}
\newcommand{\bfl}{\boldsymbol{l}}
\newcommand{\bflambda}{\mbox{\boldmath$\lambda$}}
\newcommand{\mclK}{\mathcal{K}}
\newcommand{\mclQ}{\mathcal{Q}}
\newcommand{\msfA}{\mathsf{A}}
\newcommand{\msfB}{\mathsf{B}}
\newcommand{\msfK}{\mathsf{K}}
\newcommand{\msfL}{\mathsf{L}}
\newcommand{\msfM}{\mathsf{M}}
\title{For DIRK schemes $U=B$}
\author{Anthony E. Ramirez, Abner J. Salgado}
\date{Source: arXiv:2608.27210v1 (CC BY 4.0)}
\begin{document}
\maketitle

\noindent\emph{Source and licence.} For DIRK schemes $U=B$. Version arXiv:2608.27210v1 (CC BY 4.0), distributed by arXiv under the Creative Commons Attribution 4.0 International licence, http://creativecommons.org/licenses/by/4.0/, as recorded in the arXiv licence metadata for this exact version and shown on the abstract page https://arxiv.org/abs/2608.27210v1. Full licence notice: \texttt{licenses/arxiv-2608.27210-CC-BY-4.0.txt}.\par\smallskip

\noindent\textit{Editorial scope.} Counted unit: the article's theorem theorem (source0.tex lines 596--598). The article's complete original proof of it is delivered with it (source0.tex lines 599--698, one \texttt{proof} environment of 100 lines). No mathematics is written, repaired or re-derived: the statement and the proof are the article's own lines, in the article's own notation, and no retained line is substituted or re-flowed.  No further source-local context is needed: every cross-reference of every retained line resolves inside the two delivered spans.  Nothing was omitted from any retained span, so the package carries no omission record.  No retained line contains a citation, so the wrapper carries no bibliography and no bibliography entry.  No retained line contains a figure environment, an image-inclusion command or drawing code, so no image is delivered and the package declares no asset dependency.  Editorial formatting only, in addition: an AMS \texttt{amsart} preamble that restates the article's own declarations and macros used by the retained lines, namely the environments \texttt{theorem} on separate counters, together with the standard packages those lines need.  The retained environments print in this document as Theorem 1 in order of appearance, whereas the article prints the same items with the numbering of its own full document.  The complete native source of the article as distributed in the arXiv e-print for this exact version is bundled unchanged as \texttt{sources/208696/source0.tex}.
\begin{theorem}[equivalence \AJS{for $s=3$}]
  For any $3$-stage DIRK scheme with an invertible matrix $\msfA$, $U$-stability is equivalent to $B$-stability.
\end{theorem}

\begin{proof}
  \AJS{As in the two-stage case, both notions require that, for $i=1,2,3$, $b_i>0$; so we focus on the spectra of matrices $\tilde{\msfM}$ and $\msfK$, which is to be defined}.
  
  We rewrite the quadratic form $\mclQ_3$ as
  \[
    \mclQ_3(w_1,w_2,w_3,w_4) = \mclK_3(w_2-w_1,w_3 - w_2, w_4 - w_3),
  \]
  where the coefficient matrix $\msfK$ has the following entries:
  \begin{align*}
    k_{11} &= \delta_1, \quad k_{22} = \delta_2, \quad k_{33} = \delta_3, \\
    k_{12} &= k_{21} = \frac{1}{2} \left( -\frac{\nu_{12}}{a_{11}} - \frac{\nu_{13}}{a_{11}} \right), \\
    k_{13} &= k_{31} = \frac{1}{2} \left( -\frac{\nu_{13}}{a_{11}} - \frac{\nu_{23}}{a_{22}} + \frac{a_{21}\nu_{23}}{a_{11}a_{22}} \right), \\
    k_{23} &= k_{32} = \frac{1}{2} \left( -\frac{\nu_{23}}{a_{22}} \right).
  \end{align*}

  Let us now compute the entries of $\tilde{\msfM}$. Direct computations yield:
  \begin{equation*}
    \msfA^{-\top}\msfB =
    \begin{bmatrix}
      \frac{b_1}{a_{11}} & -\frac{a_{21}b_2}{a_{11}a_{22}} & \frac{(-a_{22}a_{31}+a_{21}a_{32})b_3}{a_{11}a_{22}a_{33}} \\
      0 & \frac{b_2}{a_{22}} & -\frac{a_{32}b_3}{a_{22}a_{33}} \\
      0 & 0 & \frac{b_3}{a_{33}}
    \end{bmatrix},
    \qquad\qquad
    \msfB\msfA^{-1} =
    \begin{bmatrix}
      \frac{b_1}{a_{11}} & 0 & 0 \\
      -\frac{a_{21}b_2}{a_{11}a_{22}} & \frac{b_2}{a_{22}} & 0 \\
      \frac{(-a_{22}a_{31}+a_{21}a_{32})b_3}{a_{11}a_{22}a_{33}} & -\frac{a_{32}b_3}{a_{22}a_{33}} & \frac{b_3}{a_{33}}
    \end{bmatrix}.
  \end{equation*}
  Notice that the upper left $2 \times 2$ block is the same as in the two stage case. The row sums of $\msfA^{-\top}\msfB$ explicitly define the elements of the extrapolation vector $\bflambda$.

  We define the change-of-basis matrix
  \begin{equation*}
    \msfL \coloneqq
    \begin{bmatrix}
      1 & 0 & 0 \\
      1 & 1 & 0 \\
      1 & 1 & 1
    \end{bmatrix}
    = \begin{bmatrix}
        \bfl_1 & \bfl_2 & \bfl_3
      \end{bmatrix} \in \Real^{3 \times 3}.
  \end{equation*}
  We now compute $\frac{1}{2} \bfl_i^\top \tilde{\msfM} \bfl_i$:
  \begin{align*}
    \frac{1}{2} \bfl_3^\top \tilde{\msfM} \bfl_3 &= \frac{1}{2} \left( 2\left(\frac{b_3}{a_{33}}\right) - \lambda_3^2 \right) = \frac{1}{2}(2\lambda_3 - \lambda_3^2) = \delta_3 = k_{33}. \\
    \frac{1}{2} \bfl_2^\top \tilde{\msfM} \bfl_2 &= \frac{1}{2} \left( 2\left( \frac{b_2}{a_{22}} - \frac{a_{32}b_3}{a_{22}a_{33}} + \frac{b_3}{a_{33}} \right) - (\lambda_2+\lambda_3)^2 \right) = \frac{1}{2} \left( 2(\lambda_2+\lambda_3) - (\lambda_2+\lambda_3)^2 \right) = \delta_2 = k_{22}. \\
    \frac{1}{2} \bfl_1^\top \tilde{\msfM} \bfl_1 &= \frac{1}{2} \left( 2(\lambda_1+\lambda_2+\lambda_3) - (\lambda_1+\lambda_2+\lambda_3)^2 \right) = \delta_1 = k_{11},
  \end{align*}
  and we see that the diagonal entries of $\frac12\msfL^\top \tilde{\msfM} \msfL$ coincide with those of $\msfK$.
  
  We now evaluate the off-diagonal entries of $\frac12\msfL^\top \tilde{\msfM} \msfL$. To compute $\frac{1}{2} \bfl_2^\top \tilde{\msfM} \bfl_3$ we observe that the operation $\bfl_2^\top \msfA^{-\top} \msfB \bfl_3$ extracts the sum of the second and third rows of the third column of $\msfA^{-\top}\msfB$, which is 
  \[
    \frac{b_3}{a_{33}} - \frac{a_{32}b_3}{a_{22}a_{33}} = \lambda_3\left(1 - \frac{a_{32}}{a_{22}}\right).
  \]
  The transpose block gives $\bfl_2^\top \msfB \msfA^{-1} \bfl_3 = \frac{b_3}{a_{33}} = \lambda_3$. Gathering,
  \[
    \frac{1}{2} \bfl_2^\top \tilde{\msfM} \bfl_3 = \frac{1}{2} \left[ \lambda_3 \left(1 - \frac{a_{32}}{a_{22}}\right) + \lambda_3 - (\lambda_2 + \lambda_3)\lambda_3 \right] 
    = -\frac{1}{2a_{22}} \lambda_3 \left[ a_{32} + a_{22}(-2 + \lambda_2 + \lambda_3) \right] = -\frac{\nu_{23}}{2a_{22}} = k_{23}.
  \]

  We now compute $\frac{1}{2} \bfl_1^\top \tilde{\msfM} \bfl_2$. We have 
  \[
    \bfl_1^\top \msfA^{-\top} \msfB \bfl_2 = \lambda_2\left(1 - \frac{a_{21}}{a_{11}} \right) + \lambda_3\left(1 - \frac{a_{31}}{a_{11}} \right).
  \]
  The transpose block gives $\bfl_1^\top \msfB \msfA^{-1} \bfl_2 = \lambda_2 + \lambda_3$. We also have that
  \begin{align*}
    2 k_{12} &= \lambda_2\left(1 - \frac{a_{21}}{a_{11}}\right) + \lambda_3\left(1 - \frac{a_{31}}{a_{11}}\right) + (\lambda_2 + \lambda_3) - (\lambda_1+\lambda_2+\lambda_3)(\lambda_2+\lambda_3) \\
    &= -\frac{a_{21}}{a_{11}}\lambda_2 - \frac{a_{31}}{a_{11}}\lambda_3 + (\lambda_2+\lambda_3)\left[ 2 - \lambda_1 - \lambda_2 - \lambda_3 \right].
  \end{align*}
  To match the Bochner stability formulation, we expand $-\frac{\nu_{12}}{a_{11}} - \frac{\nu_{13}}{a_{11}}$ directly from their definitions:
  \begin{align*}
    -\frac{\nu_{12}}{a_{11}} - \frac{\nu_{13}}{a_{11}} &= -\frac{a_{21}}{a_{11}} \left[ \lambda_2(1-\lambda_3) + (2-\lambda_3)\lambda_3 + \lambda_3(-2+\lambda_2+\lambda_3) \right] \\
    &\quad - \left[ \lambda_2(-2+\lambda_1+\lambda_2) + 2(-1+\lambda_2)\lambda_3 + \lambda_3^2 \right] - \frac{a_{31}}{a_{11}}\lambda_3 - \lambda_1\lambda_3.
  \end{align*}
  The bracketed term factored by $-\frac{a_{21}}{a_{11}}$ simplifies to $\lambda_2$. The remaining terms evaluate to $(\lambda_2+\lambda_3)[2 - \lambda_1 - \lambda_2 - \lambda_3] - \frac{a_{31}}{a_{11}}\lambda_3$, thus 
  \[
    \frac{1}{2} \bfl_1^\top \tilde{\msfM} \bfl_2 = \frac{1}{2} \left( -\frac{\nu_{12}}{a_{11}} - \frac{\nu_{13}}{a_{11}} \right) = k_{12}.
  \]

  Finally, we find $\frac{1}{2} \bfl_1^\top \tilde{\msfM} \bfl_3$. The term $\bfl_1^\top \msfA^{-\top} \msfB \bfl_3$ sums the entire third column of $\msfA^{-\top} \msfB$, yielding 
  \[
    \bfl_1^\top \msfA^{-\top} \msfB \bfl_3 = \lambda_3 \left[ 1 - \frac{a_{32}}{a_{22}} - \frac{a_{31}}{a_{11}} + \frac{a_{21}a_{32}}{a_{11}a_{22}} \right].
  \]
  The transpose block gives $\lambda_3$. Therefore,
  \[
    2 k_{13} = \lambda_3 \left[ 2 - \frac{a_{31}}{a_{11}} - \frac{a_{32}}{a_{22}} + \frac{a_{21}a_{32}}{a_{11}a_{22}} \right] - (\lambda_1+\lambda_2+\lambda_3)\lambda_3 
    = \lambda_3 \left[ 2 - \lambda_1 - \lambda_2 - \lambda_3 - \frac{a_{31}}{a_{11}} - \frac{a_{32}}{a_{22}} + \frac{a_{21}a_{32}}{a_{11}a_{22}} \right].
  \]
  We expand $-\frac{\nu_{13}}{a_{11}} - \frac{\nu_{23}}{a_{22}} + \frac{a_{21}\nu_{23}}{a_{11}a_{22}}$ from their definitions to match:
  \begin{align*}
    -\frac{\nu_{13}}{a_{11}} - \frac{\nu_{23}}{a_{22}} + \frac{a_{21}\nu_{23}}{a_{11}a_{22}} &= \lambda_3 \left[ -\frac{a_{31}}{a_{11}} - \lambda_1 - \frac{a_{21}}{a_{11}}(-2+\lambda_2+\lambda_3) \right] \\
    &\quad + \lambda_3 \left[ -\frac{a_{32}}{a_{22}} + 2 - \lambda_2 - \lambda_3 \right] + \lambda_3 \left[ \frac{a_{21}}{a_{11}}\left( \frac{a_{32}}{a_{22}} - 2 + \lambda_2 + \lambda_3 \right) \right].
  \end{align*}
  Grouping the terms factored by $\frac{a_{21}}{a_{11}}$ gives us $\frac{a_{21}a_{32}}{a_{11}a_{22}}$. The remaining terms match the expansion above, hence $\frac{1}{2} \bfl_1^\top \tilde{\msfM} \bfl_3 = k_{13}$.
  
  We have proved that $\msfK = \frac{1}{2} \msfL^\top \tilde{\msfM} \msfL$, and this shows that $\msfK$ and $\frac12\msfM$ have the same spectrum. One of these matrices is positive semi-definite, if and only if the other one is.
\end{proof}

\end{document}
